% Inference cost + environmental footprint table for the 7-way baseline
% comparison. Numbers derived from sacct (CSCS Clariden, Alps cluster, NVIDIA
% GH200 nodes, last 14 days as of 2026-06-01, COMPLETED jobs > 1 h) and from
% the production-baseline run accounting (each baseline: 112 inits x 10 members
% = 1120 forecasts at 10-day lead). Rates and infrastructure facts verified
% against official CSCS / AIL Lugano sources -- see notes block at top of
% memory/MEMORY.md ("Inference cost and environmental footprint, Verified
% sources").
%
% Required packages (already in main preamble per experiments_tables.tex):
%   amsmath, booktabs, xcolor[table], array
%
% Reference via \ref{tab:inference-cost}.
%
% Per-member wall-clock formula: (mean job time in min) / (14 init_times) /
% (10 members) for week-helper jobs. Trained-prob baselines (aifsens, atlas,
% fcn3) here use fewer sample jobs in the 14-day window because their
% production runs completed earlier; the per-job mean is still representative.
%
% Total NH per production-baseline run = 8 week-helper jobs x ~3-4 h x 1 node.
% CHF rate: Swiss academia tier 3 (> 10000 NH/yr) at CHF 3.84/NH, with the
% ETH 25% discount applied: CHF 2.88 / node-hour.
% kWh per NH at GH200 node: peak ~5.6 kW per node (4x Grace-Hopper at ~1 kW
% each plus interconnect), conservative average ~4.0 kW under steady inference
% load, x PUE 1.2 -> ~4.8 kWh per node-hour delivered.

\begin{table}[t]
  \centering
  \small
  \setlength{\tabcolsep}{4pt}
  \renewcommand{\arraystretch}{1.20}
  \caption{Per-baseline inference cost on CSCS Alps (NVIDIA GH200) for the
           7-way production-baseline study (112 inits $\times$ 10 members =
           1120 forecasts per baseline, 10-day lead). \textbf{Pure inference
           time per member} (sec/mbr) is the mean per-call duration between
           "Inference starting!" and "Inference complete" markers in the
           earth2studio logs (median for SFNO, where the mean is inflated by
           SIGSEGV-recovery retries); for Atlas, derived from the tqdm
           progress-bar total time (61 rollout steps $\times$ 37.5 s/step).
           This excludes data fetch (ARCO/CDS), model load, container
           startup, and zarr write. NH = compute node-hours (1 Alps node =
           4 Grace-Hopper superchips). Cost uses the CSCS Swiss-academia
           rate at tier 3 (> 10\,000 NH/year, CHF 3.84/NH) with the
           ETH-Zurich 25\% discount: CHF 2.88/NH delivered. Energy assumes
           4.0 kW node power under steady inference and PUE 1.2 (peak rated
           5.6 kW; \texttt{<} 1.2 PUE publicly disclosed by CSCS, annualised
           figure not disclosed). The data centre runs on 100\% Ticino
           hydropower (Aziende Industriali di Lugano contract, 36 GWh/yr;
           cscs.ch news 2021); cooling water is drawn from Lake Lugano at
           45 m depth (inlet 6\,$^{\circ}$C) and returned at $\sim$12 m depth (CSCS
           lake-water fact sheet); waste heat is delivered to the CSCS
           office and USI/SUPSI campus year-round (\texttt{~}742 MWh
           thermal/yr), with a 6 MW (scalable to 30 MW) AIL district-heating
           plant under construction. The fastest baseline is SFNO modes10
           (59 s/member, 8 spherical-operator blocks); AIFS-CRPS is second
           (75 s/member, single forward pass with conditional-LN noise).
           Atlas is \emph{the dramatic outlier} at 38 minutes per member -
           roughly 25-40x slower than the other six - because its
           stochastic-interpolant Diffusion Transformer requires multiple
           denoising steps per rollout step. Our post-hoc perturbation
           baselines (59-125 s/member) are competitive with the
           single-forward-pass trained-probabilistic baselines and
           \emph{vastly} cheaper than diffusion-style stochastic sampling.}
  \label{tab:inference-cost}
  \begin{tabular}{@{}l l rrrrr@{}}
    \toprule
    \textbf{Role}
        & \textbf{Model}
        & \textbf{s/mbr$^{\dag}$}
        & \textbf{CS h$^{\S}$}
        & \textbf{NH$^{\ddag}$}
        & \textbf{Cost (CHF)$^{\ddag}$}
        & \textbf{kWh$^{\ddag}$} \\
    \midrule

    \rowcolor{phaseBg}
    \multicolumn{7}{@{}l}{\textbf{Post-hoc weight perturbation} (this work)} \\
    & \mbadge{sfnoC}{sfno\_modes10}      & \textbf{59}  & 0.20  & 15.3 & 44.1  & 73.4  \\
    \rowcolor{altRow}
    & \mbadge{auroraC}{aurora\_encoder}  & 90           & 0.31  & 23.0 & 66.2  & 110.4 \\
    & \mbadge{graphcastC}{graphcast\_all} & 125         & 0.43  & 28.7 & 82.7  & 137.8 \\
    \rowcolor{altRow}
    & \mbadge{aifsC}{aifs\_perturbed}    & (pending)$^{\star}$    & (pending)  & (pending)   & (pending) & (pending)  \\
    \midrule

    \rowcolor{phaseBg}
    \multicolumn{7}{@{}l}{\textbf{Trained probabilistic}} \\
    & \mbadge{aifsensC}{aifsens}         & 75            & 0.26  & 18.3 & 52.7  & 87.8  \\
    \rowcolor{altRow}
    & \mbadge{fcn3C}{fcn3}               & 114           & 0.40  & 26.6 & 76.6  & 127.7 \\
    & \mbadge{atlasC}{atlas}             & \textbf{2280} & \textbf{7.92}  & 43.5 & 125.3 & 208.8 \\
    \midrule

    \rowcolor{phaseBg}
    \multicolumn{7}{@{}l}{\textbf{Classical reference} (off-site)} \\
    & \mbadge{ifsC}{ifs\_ens}          & \multicolumn{5}{l}{ECMWF MARS download only; no CSCS compute} \\

    \bottomrule
  \end{tabular}
  \vspace{2pt}

  \scriptsize
  $^{\dag}$\,Pure inference seconds per ensemble member per 10-day forecast
  (excludes data fetch, model load, container startup, zarr write). Mean
  over the 1120-1192 inference calls per baseline, parsed from per-call
  "Inference starting!"/"Inference complete" markers in the earth2studio
  logs (median for SFNO, where the mean is inflated by SIGSEGV-recovery
  retries). Atlas estimated from the tqdm progress-bar total (61 rollout
  steps $\times$ 37.5 s/step).
  $^{\S}$\,\textbf{CS h} = wall-clock hours for a 50-member case-study
  forecast on a single Alps node (4 GH200 in parallel):
  $50 \times \text{s/mbr} / 4 / 3600$. Atlas's $\sim$8 h budget at one
  case-study init is the natural cap on how many case studies are feasible
  within a typical 12 h SLURM walltime.
  $^{\star}$\,\textbf{aifs\_perturbed} = the best post-hoc weight
  perturbation of the AIFS v1 deterministic checkpoint, identified from
  the in-flight Phase 1 and Phase 2 AIFS ablations (winner pending eval).
  Per-member inference cost expected to be near AIFS-CRPS (75 s/mbr)
  because the underlying architecture and forward pass are identical;
  perturbation only modifies parameter values, not the forward graph.
  $^{\ddag}$\,Per production-baseline run (1120 forecasts).
\end{table}
